% auto-generated by make exhibits -- do not edit
\begin{table}[!htbp]\centering\footnotesize
\caption{What each type of government borrows for}\label{tab:d2}
\begin{tabular}{@{}lc>{\raggedright\arraybackslash}p{0.55\linewidth}@{}}
\toprule
Government type & Project \$B & Largest purposes (share of the type's project dollars) \\
\midrule
School districts & 356 & education 98\%; general government 1\%; parks / recreation / culture 0\% \\
Municipalities & 386 & general government 36\%; water / sewer / environment 17\%; education 12\% \\
Counties & 120 & education 25\%; transportation 18\%; public safety / justice 16\% \\
Townships & 31 & education 30\%; water / sewer / environment 27\%; general government 15\% \\
Special districts & 85 & water / sewer / environment 33\%; transportation 19\%; utilities 19\% \\
\bottomrule
\end{tabular}
\begin{minipage}{0.96\linewidth}\vspace{2pt}\scriptsize
\emph{Notes:} Same universe as Table \ref{tab:d1}; project dollars by the document's primary major function, excluding pure financing mechanics. Shares are of the type's own classified project dollars.
\end{minipage}
% BEGIN READING (strip for submission)
\begin{minipage}{0.96\linewidth}\vspace{4pt}\scriptsize
\textbf{Reading.} Each type of government is a different bundle of services, which is why the consent requirement cannot be one institution: the goods behind the ballot differ by who is asking.
\end{minipage}
% END READING
\end{table}
